\subsection{半单模的反模}

设 A 是一个环。把单 A-模的类组成的集合记作 S。对每个 \(\lambda \in S\)，选取一个单 A-模 \(S_{\lambda}\)，其类为 \(\lambda\)，并把 \(S_{\lambda}\) 的自同态体之反环记作 \(D_{\lambda}\)。我们把 \(S_{\lambda}\) 视为 \((A, D_{\lambda})\)-双模。

设 \(M\) 是一个半单 A-模，\(B\) 是 \(M\) 的自同态环。把 \(M\) 的双交换子记作 \(C\)。对每个 \(\lambda \in \mathcal{S}\)，把左 \((D_{\lambda}, B)\)-双模 \(\mathrm{Hom}_A(S_{\lambda}, M)\) 记作 \(V_{\lambda}\)。最后，把 A-模 \(M\) 的支集记作 \(\mathcal{S}_M\)（VIII, 第 66 页）；它也是元素 \(\lambda\)（属于 \(\mathcal{S}\)）组成的集合，这些元素使 \(V_{\lambda}\) 非零。

注 1. —— A-模 M 的典范描述 \(\alpha_{M}\) 是左 (A, B)-双模的同构。由 VIII, 第 71 页命题 9 的推论，映射 \(f \mapsto (\operatorname{Hom}(1_{\mathrm{S}_{\lambda}}, f))_{\lambda \in \mathcal{S}_{\mathrm{M}}}\)（从 B 到 \(\prod_{\lambda \in \mathcal{S}_{\mathrm{M}}} \operatorname{End}_{\mathrm{D}_{\lambda}}(\mathrm{V}_{\lambda})\)）是环同构。

命题 5. —— a) M 的反模是半单的。

\begin{enumerate}
\def\labelenumi{\alph{enumi})}
\setcounter{enumi}{1}
\item
  对每个 \(\lambda \in \mathcal{S}_{\mathrm{M}}\)，B-模 \(V_{\lambda}\) 是单模，其交换子等于 \((D_{\lambda})_{V_{\lambda}}\)。
\item
  映射 \(\lambda \mapsto \mathrm{cl}(\mathrm{V}_{\lambda})\) 是从 A-模 M 的支集到其反模的支集的双射。
\item
  对每个 \(\lambda \in \mathcal{S}_{\mathrm{M}}\)，B-子模 \(\mathrm{M}_{\lambda}\) 是 B-模 M 的 \(\mathrm{V}_{\lambda}\) 型同型分量，且 \(\mathrm{V}_{\lambda}\) 在 M 中的重数等于 \(\dim_{\mathrm{D}_{\lambda}}(\mathrm{S}_{\lambda})\)。
\item
  对 \(s \in S\)，把映射 \(\varphi \mapsto \varphi(s)\)（从 \(V_{\lambda} = \operatorname{Hom}_{\mathrm{A}}(\mathrm{S}_{\lambda}, \mathrm{M})\) 到 M）记作 \(\widetilde{s}\)。它是 B-线性的。所得到的映射 \(s \mapsto \widetilde{s}\)（从 \(S_{\lambda}\) 到 \(\operatorname{Hom}_{\mathrm{B}}(\mathrm{V}_{\lambda}, \mathrm{M})\)）是 \((\mathrm{A}, \mathrm{D}_{\lambda})\)-双模的同构。
\end{enumerate}

设 \(\lambda \in \mathcal{S}_{\mathrm{M}}\)。把环 \(\mathrm{End}_{\mathrm{D}_{\lambda}}(\mathrm{V}_{\lambda})\) 记作 \(\mathrm{E}_{\lambda}\)；由于 \(\mathrm{V}_{\lambda}\) 是非零 \(\mathrm{D}_{\lambda}\)-向量空间，它是单 \(\mathrm{E}_{\lambda}\)-模（VIII, 第 45 页，例 3），其交换子等于 \((\mathrm{D}_{\lambda})_{\mathrm{V}_{\lambda}}\)（VIII, 第 82 页，定理 2 的推论 1）。由于 \(\mathrm{E}_{\lambda}\) 是 B-模 \(\mathrm{V}_{\lambda}\) 的位似环（VIII, 第 71 页，命题 9 的推论），这就证明了 b)。

M 的典范描述 \(\alpha_{M}\) 定义了一个同构 \(\alpha_{\lambda}\)（从 \(V_{\lambda}\otimes_{D_{\lambda}^{\circ}}S_{\lambda}\) 到 \(M_{\lambda}\)）。由于 \(V_{\lambda}\) 是单 B-模，B-模 \(V_{\lambda}\otimes_{D_{\lambda}^{\circ}}S_{\lambda}\) 是 \(V_{\lambda}\) 型同型模（VIII, 第 61 页，命题 1）；因此 B-模 \(M_{\lambda}\) 亦然，这就证明了 a)。

由上面的注 1，存在 B 的元素 \(e_{\lambda}\)（\(\lambda\) 取遍 \(\mathcal{S}_{\mathrm{M}}\)），使得 \((e_{\lambda})_{\mathrm{V}_{\lambda}} = 1_{\mathrm{V}_{\lambda}}\)，且 \((e_{\lambda})_{\mathrm{V}_{\mu}} = 0\)（当 \(\mu \in \mathcal{S}_{\mathrm{M}}\) 而 \(\mu \neq \lambda\)）。因此单 B-模 \(V_{\lambda}\) 两两不同构，这就证明了 c) 及 d) 的第一条断言。B-模 \(M_{\lambda}\) 同构于 \(V_{\lambda} \otimes_{D_{\lambda}^{\circ}} S_{\lambda}\)，故 \(\dim_{D_{\lambda}}(S_{\lambda})\) 是 \(V_{\lambda}\) 在 M 中的重数（II, §3, No.~7, 第 255 页，推论 1）。

映射 \(\sum_{\lambda\in\mathcal{S}_{M}}\alpha_{\lambda}\)（从 \(\bigoplus_{\lambda\in\mathcal{S}_{M}}V_{\lambda}\otimes_{D_{\lambda}^{\circ}}S_{\lambda}\) 到 M）给出半单 B-模 M 的一个描述（VIII, 第 69 页，定义 5）。由 VIII, 第 70 页命题 8, b)，对每个 \(\lambda\in\mathcal{S}_{M}\)，e) 中所述的从 \(S_{\lambda}\) 到 \(\mathrm{Hom}_{\mathrm{B}}(\mathrm{V}_{\lambda},\mathrm{M})\) 的映射是双射且 \(D_{\lambda}\)-线性。由于它显然是 A-线性的，这就证明了 e)。

注 2. —— 由证明可知，映射

\[
\sum_ {\lambda \in \mathcal {S} _ {\mathrm{M}}} V _ {\lambda} \otimes_ {D _ {\lambda} ^ {\mathrm{o}}} S _ {\lambda} \longrightarrow M
\]

（由 M 的典范描述所诱导）是 M 的反模的一个描述。

命题 6. —— a) 视为 (A, B°)-双模时，M 是半单的。

\begin{enumerate}
\def\labelenumi{\alph{enumi})}
\setcounter{enumi}{1}
\item
  对每个 \(\lambda \in \mathcal{S}_{\mathrm{M}}\)，\(\mathrm{M}_{\lambda}\) 是 M 的单 (A, B°)-子双模。
\item
  对 M 的每个 \((\mathrm{A},\mathrm{B}^{\circ})\)-子双模 N，存在一个唯一的子集 \(\Lambda\)（\(\mathcal{S}_{\mathrm{M}}\) 的），使得 N 等于 \(\oplus_{\lambda \in \Lambda}\mathrm{M}_{\lambda}\)。
\end{enumerate}

设 \(\lambda\) 属于 \(\mathcal{S}_{\mathrm{M}}\)。左 A-模 \(S_{\lambda}\) 与右 B-模 \(V_{\lambda}\) 都是单模，且 \(D_{\lambda}\) 是 \(S_{\lambda}\) 的自同态体之反环。由 VIII, 第 63 页的推论 2，(A, B°)-双模 \(S_{\lambda} \otimes_{D_{\lambda}} V_{\lambda}\) 是单的，而与它同构的 \(M_{\lambda}\) 亦然。这就证明了 b)，a) 随之可得。

若 \(\lambda\) 与 \(\mu\) 在 \(S_{M}\) 中不同，则 \(M_{\lambda}\) 与 \(M_{\mu}\) 作为 A-模不同构，更不必说作为 \((A,B^{\circ})\)-双模了。断言 c) 由 VIII, 第 68 页的推论 2 得出。

命题 7. —— a) 对 M 的双交换子 C 的每个元素 c 及每个 \(\lambda\in\mathcal{S}_{M}\)，存在唯一的元素 \(c_{\lambda}\)（属于 \(\operatorname{End}_{\mathrm{D}_{\lambda}}(\mathrm{S}_{\lambda})\)），使得对每个 \(\varphi\in\operatorname{Hom}_{\mathrm{A}}(\mathrm{S}_{\lambda},\mathrm{M})\) 及每个 \(s\in S_{\lambda}\)，我们有 \(c\varphi(s)=\varphi(c_{\lambda}s)\)。

\begin{enumerate}
\def\labelenumi{\alph{enumi})}
\setcounter{enumi}{1}
\item
  给 \(S_{\lambda}\)（\(\lambda\) 取遍 \(S_{M}\)）赋以由 a) 定义的 C-模结构。则典范映射 \(\alpha_{M}\)（从 \(\bigoplus_{\lambda\in S_{M}}S_{\lambda}\otimes_{D_{\lambda}}V_{\lambda}\) 到 M）是 \((\mathrm{C},\mathrm{B}^{\circ})\)-双模的同构。
\item
  映射 \(c \mapsto (c_{\lambda})_{\lambda \in \mathcal{S}_{\mathrm{M}}}\) 是从 C 到 \(\prod_{\lambda \in \mathcal{S}_{\mathrm{M}}}\mathrm{End}_{\mathrm{D}_{\lambda}}(\mathrm{S}_{\lambda})\) 的同构。
\end{enumerate}

断言 a) 与 c) 由 VIII, 第 71 页命题 9 得出，因为典范映射 \(\alpha_{\mathrm{M}}\)（从 \(\bigoplus_{\lambda \in \mathcal{S}} (\mathrm{S}_{\lambda} \otimes_{\mathrm{D}_{\lambda}} \mathrm{W}_{\lambda})\) 到 M）给出 B-模 M 的一个描述（注 2）。此外，\(\alpha_{\mathrm{M}}\) 是 (C, B°)-线性的，这就证明了 b)。

注 3. —— 给 \(S_{\lambda}\)（对每个 \(\lambda \in S_{M}\)）赋以命题 7, a) 给出的 C-模结构。若在命题 5（VIII, 第 85 页）中把 A 换成 B、把 B 换成 C，则可看出：对每个 \(\lambda \in S_{M}\)，左 C-模 \(S_{\lambda}\) 是单模，其交换子为 \(D_{\lambda}\)；C-模 M 的 \(S_{\lambda}\) 型同型分量等于 \(M_{\lambda}\)；且映射 \(\lambda \mapsto \mathrm{cl}_{\mathrm{C}}(\mathrm{S}_{\lambda})\) 是从 A-模 M 的支集到 C-模 M 的支集的双射。最后注意，从 \(S_{\lambda}\) 到 M 的 A-线性映射与 C-线性映射相同，M 的 A-子模与 C-子模相同，且环 \(\operatorname{End}_{\mathrm{A}}(\mathrm{M})\) 与 \(\operatorname{End}_{\mathrm{C}}(\mathrm{M})\) 相等。

设 M 是一个半单 A-模。把 A-模 M 的双交换子 C 的中心记作 Z；它也是 M 的交换子 B 的中心。给 M 及 \(S_{\lambda}\)（\(\lambda \in S_{M}\)）赋以由 C-模结构经标量限制得到的 Z-模结构。对每个 \(\lambda \in S_{M}\)，把体 \(D_{\lambda}\) 的中心记作 \(Z_{\lambda}\)。

命题 8. —— a) 映射 \(z \mapsto (z_{\mathrm{S}_{\lambda}})_{\lambda \in \mathcal{S}_{\mathrm{M}}}\) 是从 Z 到诸体 \(\mathbf{Z}_{\lambda}\) 的乘积的同构。

\begin{enumerate}
\def\labelenumi{\alph{enumi})}
\setcounter{enumi}{1}
\item
  A-模 M 是同型的且非零，当且仅当 Z 是一个体。
\item
  设 \(\Lambda\) 是 \(\mathcal{S}_{\mathrm{M}}\) 的一个子集。用 \(e_{\Lambda}\) 表示 \(\mathrm{Z}\) 中唯一的元素，它满足：\((e_{\Lambda})_{\mathrm{S}_{\lambda}} = 1_{\mathrm{S}_{\lambda}}\)（对 \(\lambda \in \Lambda\)），且 \((e_{\Lambda})_{\mathrm{S}_{\lambda}} = 0\)（对 \(\lambda \in \mathcal{S}_{\mathrm{M}} - \Lambda\)）。我们有 \((e_{\Lambda})_{\mathrm{M}_{\lambda}} = 1_{\mathrm{M}_{\lambda}}\)（对 \(\lambda \in \Lambda\)）及 \((e_{\Lambda})_{\mathrm{M}_{\lambda}} = 0\)（对 \(\lambda \in \mathcal{S}_{\mathrm{M}} - \Lambda\)）。
\item
  若 M 的支集 \(S_{M}\) 有限，则映射 \(\Lambda \mapsto e_{\Lambda}Z\) 是从 \(S_{M}\) 的子集组成的集合到 Z 的理想组成的集合的双射，而映射 \(a \mapsto aM\) 是从 Z 的理想组成的集合到 M 的 \((A,B^{o})\)-子双模组成的集合的双射。这些双射都是有序集的同构。逆双射把 M 的 \((A,B^{o})\)-子双模 N 映到由 Z 中把 N 送入 M 的元素 z 组成的理想。
\end{enumerate}

对 \(\lambda \in \mathcal{S}_{\mathrm{M}}\)，\(Z_{\lambda}\) 是交换子 \(D_{\lambda}\) 与双交换子 \(C_{\lambda}\)（A-模 \(S_{\lambda}\) 的）的公共中心。由上面的命题 7, c)，映射 \(c \mapsto (c_{\lambda})_{\lambda \in \mathcal{S}_{\mathrm{M}}}\) 是从 C 到 \(\prod_{\lambda \in \mathcal{S}_{\mathrm{M}}} C_{\lambda}\) 的同构。限制到中心上，我们得到同构 \(z \mapsto (z_{S_{\lambda}})_{\lambda \in \mathcal{S}_{\mathrm{M}}}\)（从 Z 到 \(\prod_{\lambda \in \mathcal{S}_{\mathrm{M}}} Z_{\lambda}\)），由此得 a)。

环 \(\prod_{\lambda \in \mathcal{S}_{\mathrm{M}}}Z_{\lambda}\) 是一个体，当且仅当集合 \(\mathcal{S}_{\mathrm{M}}\) 只有一个元素，由此得 b)。

断言 c) 由命题 7, a) 得出。设 \(\mathcal{S}_{\mathrm{M}}\) 有限。由 a) 及 I, §8, No.~10, 第 109 页的命题 8，映射 \(\Lambda \mapsto e_{\Lambda}\mathrm{Z}\) 是从 \(\mathfrak{P}(\mathcal{S}_{\mathrm{M}})\) 到 Z 的理想组成的集合的有序集同构。设 \(\Lambda\) 是 \(\mathcal{S}_{\mathrm{M}}\) 的一个子集。由 c)，我们有关系式 \(e_{\Lambda}\mathrm{ZM} = e_{\Lambda}\mathrm{M} = \bigoplus_{\lambda \in \Lambda}\mathrm{M}_{\lambda}\)；由 VIII, 第 86 页的命题 6, c)，剩下只需描述逆双射。而 \(z \in \mathrm{Z}\) 把 M 送入 \(\bigoplus_{\lambda \in \Lambda}\mathrm{M}_{\lambda}\) 当且仅当 \(z = e_{\Lambda}z\)，即 \(z \in e_{\Lambda}\mathrm{Z}\)。

推论. —— 假设 A 是代数闭交换域（域）K 上的代数，且 M 是一个作为 K 上向量空间有限维的半单 A-模。对每个元素 \(\lambda\)（在 \(S_{M}\) 中），用 \(e_{\lambda}\) 表示 M 的投影算子，其像为 \(M_{\lambda}\)，核为 \(\oplus_{\lambda\neq\mu}M_{\mu}\)。则 \((e_{\lambda})_{\lambda\in\mathcal{S}_{M}}\) 是 K 上向量空间 Z 的一个基。

由于 M 是 K 上有限维向量空间，且是非零子模族 \((\mathrm{M}_{\lambda})_{\lambda\in\mathcal{S}_{\mathrm{M}}}\) 的直和，故集合 \(S_{M}\) 有限，且每个空间 \(S_{\lambda}\)（\(\lambda\in S_{M}\)）都在 K 上有限维。由于体 K 代数闭，我们有 \(D_{\lambda}=Z_{\lambda}=K\)（VIII, 第 47 页，定理 1），且映射 \(z\mapsto(z_{\mathrm{S}_{\lambda}})_{\lambda\in\mathcal{S}_{\mathrm{M}}}\) 是从 Z 到 \(K^{S_{M}}\) 的同构（命题 8, a)）。于是本推论由命题 8 的 c) 得出。

